\centerline{\bf \Huge Kомбигеомная солянка.}

\begin{flushright}
        \scriptsize{Страшный сон Яны и Саши}
\end{flushright}

\problem{0} 
На прямой выбраны несколько отрезков.

a) \textbf{Теорема Хелли на прямой.} Известно, что любые два отрезка пересекаются. Докажите, что можно отметить точку, принадлежающую всем отрезкам.

б) \textbf{Цветная теорема Хелли на прямой.}  Известно, что среди любых $k + 1$ отрезка найдутся два пересекающихся. Докажите, что можно отметить на прямой $k$ красных точек так, чтобы любой из отрезков содержал какую-то красную точку.

\problem{1}
На плоскости дано конечное множество многоульников, каждые два из которых имеют общую точку. Докажите, что существует прямая, имеющая общие точки со всеми этими многоугольниками.

\problem{2}
На плоскости выбраны несколько прямоугольников со сторонами, параллельными осям координат. Известно, что любые два таких прямоугольника пересекатся. Докажите, что все такие прямоугольники пересекаются в одной точке.

\problem{3} Докажите, что теорема Хелли на прямой не верна в случае бесконечного количества множеств.

\textbf{Определение.} \textit{Хроматическое число графа — минимальное число цветов, в которые можно раскрасить вершины графа так, чтобы концы любого ребра имели разные цвета.}

\problem{4} 
Найдите хроматическое число прямой.

\problem{5} 
Докажите, что

a) Хроматическое число плоскости не меньше 3.

б) Хроматическое число плоскости не меньше 4.

в$)^*$ Хроматическое число плоскости не больше 9.

г$)^{**}$ Хроматическое число плоскости не больше 7.

\problem{6}
На плоскости дано $n$ точек и отмечены середины всех отрезков с концами в этих точках. Докажите, что различных отмеченных точек не менее $2n - 3.$